\section{The Universal Cause Is Not an Abstraction}

Our minds reach a universal concept by leaving particulars behind. We encounter this oak and that oak, attend to what belongs to both, and think \emph{tree} without thinking either trunk. The concept becomes more general as it includes more instances and says less about each. If God is imagined as the widest concept of all---being with every contour removed---then concrete action will look like a descent from purity. The safest god will be the one of whom almost nothing can be said because he does almost nothing in particular.

But a universal concept and a universal cause are nearly opposites. The concept \emph{tree} does not grow a leaf. It has no knowledge of oaks and gives them no existence. A cause, precisely as cause, is known through the actuality it gives. To say that God is the universal cause is not to say that he is a very inclusive thought. It is to say that every creature's act of being depends upon him.\endnote{Aquinas, ST I, q.~44, a.~1, treats God as the first cause of the being of things; ST I, q.~104, a.~1, treats every creature's present conservation in being. CCC 301 gives the authoritative catechetical synthesis that God upholds and sustains creation ``at every moment''; \sourceurl{https://www.vatican.va/content/catechism/en/part_one/section_two/chapter_one/article_1/paragraph_4_the_creator.html}{official English}.} The universality lies in the reach of his causality, not in a lack of reality.

Aquinas therefore denies that God belongs to a genus. Whatever belongs to a genus receives the genus under a limiting difference: animal becomes rational animal, color becomes red. God does not receive being under a contracted mode as creatures do; he is not one case of a more comprehensive reality. The same Aquinas who insists upon this transcendence also insists that God knows singular things. God causes not merely what creatures share but the whole actuality by which each exists as this creature. Divine knowledge is not forced to choose between the universal and the singular as ours often is.\endnote{Aquinas, \emph{Summa theologiae} (ST) I, q.~3, a.~5, argues that God is not in a genus; ST I, q.~14, a.~11, argues that God knows singulars because the perfections divided among lower knowers exist simply and unitedly in God. Latin loci were checked through the \sourceurl{https://www.corpusthomisticum.org/iopera.html}{Corpus Thomisticum}; the \sourceurl{https://www.newadvent.org/summa/1003.htm}{ST I, q.~3} and \sourceurl{https://www.newadvent.org/summa/1014.htm}{ST I, q.~14} pages served as working public English aids. Work, question, article, and reply govern over either site's phrasing.}

Transcendence is consequently not divine vagueness. Apophatic theology is right to deny that any created concept comprehends the divine essence. Every true likeness between Creator and creature is exceeded by an even greater unlikeness. Yet incomprehensibility does not mean that God is indeterminate in himself, and it does not revoke his freedom to make himself known. The failure to contain God in a concept is not permission to replace him with an abstraction.\endnote{See CCC 39--43 on analogical knowledge, the purification of language, and the Creator's transcendence beyond every creaturely likeness; \sourceurl{https://www.vatican.va/content/catechism/en/part_one/section_one/chapter_one/iv_how_can_we_speak_about_god.html}{official English}. This article supports natural theology and apophatic discipline. It opposes only the inference that transcendence makes personal or historical revelation intrinsically unworthy of God.}

Nor does particular providence turn the world into a sequence of arbitrary interruptions. Christian teaching joins God's ``concrete and immediate'' care for all things to the genuine causality of creatures. Rain still falls through weather, bread through earth and labor, healing through the body and medicine, and human deeds through free agents. God gives creatures not only existence but the dignity of acting. Some evils are permitted rather than positively willed; suffering is not thereby decoded as a private sentence imposed upon its victim.\endnote{CCC 302--308 treats providence, secondary causality, free cooperation, and the permission of evil; \sourceurl{https://www.vatican.va/content/catechism/en/part_one/section_two/chapter_one/article_1/paragraph_4_the_creator.html}{official English}. Aquinas, ST I, q.~22, aa.~2--3, supplies the adopted account of providence extending to all things and governing through created causes. The present claim neither assigns a special divine message to each event nor denies the ordinary work of natural, personal, medical, and social causes.}

The universal cause is therefore not kept universal by remaining idle. He is universal because nothing real falls outside the dependence by which it is, while each created nature and each secondary cause remains genuinely its own. If such a God speaks, his word need not become less divine by being addressed. An address that is nowhere, in no language, to no one, is not a more universal address. It is no address at all.
